% ============================================================
% Come si usano, argomento 21: stima bayesiana, MMSE e MAP
% ============================================================
\argomento{Stima bayesiana: MMSE e MAP}{Formule a p.~\pageref*{f:aposteriori}}

\begin{scelta}{Come scegliere la formula}
\renewcommand{\arraystretch}{1.25}
\begin{tabular}{@{}P{0.55\linewidth}P{0.39\linewidth}@{}}
  \toprule
  \textbf{Se il testo\ldots} & \textbf{si usa\ldots} \\
  \midrule
  dà una densità a priori del parametro e chiede che cosa se ne sa dopo i dati & la densità a posteriori
    \vv{aposteriori} \\
  chiede lo stimatore a minimo errore quadratico medio & \vv{mmse} \\
  chiede il valore più probabile del parametro, o il massimo a posteriori & \vv{stimamap} \\
  dà una funzione di costo diversa & lo stimatore ottimo \vv{stimatoreottimo} \\
  parla di una probabilità di successo non nota, uniforme a priori & \vv{bernoullicomposito}
    \vv{prestazionistime} \\
  chiede se lo stimatore migliora con tanti dati & \vv{consistenza} \\
  \bottomrule
\end{tabular}

MMSE e MAP guardano la stessa densità a posteriori: il primo ne prende il baricentro, il secondo il picco. Per il MAP
non serve nemmeno il denominatore della legge di Bayes.
\end{scelta}

\scheda{aposteriori}{
  \usoquando{Il parametro è aleatorio, con una densità a priori nota, e si vuole aggiornare quello che se ne sa dopo
    aver visto i dati.}
  \usocome{Numeratore: verosimiglianza per densità a priori. Denominatore: lo stesso numeratore integrato in
    $\theta$, che serve solo a far tornare l'area pari a 1.}
  \usoesempio{Moneta con $B$ uniforme, 2 teste in 3 lanci: numeratore $\beta^2 (1 - \beta)$, il cui integrale è
    $\frac{2!\, 1!}{4!} = \frac{1}{12}$ (\vv{eulero}); quindi $f_{B \mid \vett{X}}(\beta \mid \vett{x}) = 12\, \beta^2
    (1 - \beta)$.}
}

\scheda{stimatoreottimo}{
  \usoquando{Si deve progettare uno stimatore dicendo quanto costano gli errori di stima.}
  \usocome{Per ogni $\vett{x}$ si rende minimo il costo medio calcolato con la densità a posteriori: il costo
    quadratico dà l'MMSE, il costo a tacca il MAP.}
  \usoesempio{A posteriori uniforme su $[0, 2]$, costo quadratico: il costo medio è $\frac{1}{2} \int_0^2 (\hat\theta -
    \theta)^2 \dd\theta = \hat\theta^2 - 2\hat\theta + \frac{4}{3}$, minimo in $\hat\theta = 1$, la media.}
}

\scheda{mmse}{
  \usoquando{Si vuole l'errore quadratico medio più piccolo possibile: è lo stimatore bayesiano più usato.}
  \usocome{Si calcola la media della densità a posteriori.}
  \usoesempio{Moneta, $f_{B \mid \vett{X}} = 12\, \beta^2 (1 - \beta)$ (2 teste in 3 lanci): $\hat\beta_{\text{MMSE}} =
    \int_0^1 12\, \beta^3 (1 - \beta) \dd\beta = 12 \cdot \frac{1}{20} = \frac{3}{5}$, come $\frac{w + 1}{n + 2}$.}
}

\scheda{stimamap}{
  \usoquando{Si cerca il valore più plausibile del parametro, dati i dati: il picco della densità a posteriori.}
  \usocome{Si toglie il denominatore e i fattori costanti, si prende il logaritmo e si annulla la derivata rispetto a
    $\theta$ (\vv{argminmax}).}
  \usoattenzione{Se il massimo cade in un estremo dell'intervallo la derivata non si annulla: si guardano anche gli
    estremi.}
  \usoesempio{Moneta: $\frac{\dd}{\dd\beta} \left[ w \ln\beta + (n - w) \ln(1 - \beta) \right] = \frac{w}{\beta} -
    \frac{n - w}{1 - \beta} = 0$ dà $\hat\beta = \frac{w}{n}$. Con 3 teste in 3 lanci il massimo è nell'estremo
    $\beta = 1$.}
}

\scheda{bernoullicomposito}{
  \usoquando{Una probabilità di successo non nota, uniforme a priori, e $n$ prove di cui $w$ riuscite.}
  \usocome{Il MAP è la frequenza; l'MMSE è come se ai lanci veri si aggiungessero una testa e una croce immaginarie.}
  \usoesempio{3 teste in 3 lanci: il MAP dice 1 («la moneta dà sempre testa»), l'MMSE, più prudente, $\frac{4}{5}$. È
    la regola di successione di Laplace: dopo $n$ albe, la probabilità che il sole sorga domani è $\frac{n + 1}{n + 2}$.}
}

\scheda{prestazionistime}{
  \usoquando{Si confrontano i due stimatori della moneta, o si chiede quanto sbagliano in media.}
  \usocome{Prima si media sapendo $B = \beta$, quando $w$ è binomiale $\mathcal{B}(n, \beta)$; poi si media su $B$
    uniforme (\vv{mediacondc}).}
  \usoesempio{$n = 10$: errore MMSE $\frac{1}{72} \simeq 0{,}0139$, errore MAP $\frac{1}{60} \simeq 0{,}0167$. Con
    $n = 1$ sono $0{,}056$ e $0{,}167$: la differenza conta solo con pochi dati.}
}

\scheda{consistenza}{
  \usoquando{Si chiede se uno stimatore si avvicina sempre di più al parametro quando i dati crescono.}
  \usocome{Si calcola l'errore quadratico medio in funzione di $n$: se tende a zero, lo stimatore è consistente in
    media quadratica, e quindi anche in senso debole.}
  \usoesempio{Moneta: $\frac{1}{6(n + 2)}$ e $\frac{1}{6n}$ tendono a zero, quindi entrambi gli stimatori sono
    consistenti.}
}
